\documentclass[aspectratio=169,11pt]{beamer}

\usepackage[utf8]{inputenc}
\usepackage[T1]{fontenc}
\usepackage[main=ngerman,provide=*]{babel}
\usepackage{lmodern}
\usepackage{booktabs}
\usepackage{array}
\usepackage{tabularx}
\usepackage{listings}
\usepackage[most]{tcolorbox}
\usepackage{hyperref}

\hypersetup{colorlinks=true,urlcolor=black,linkcolor=black}

\definecolor{paper}{RGB}{248,247,243}
\definecolor{ink}{RGB}{26,32,44}
\definecolor{muted}{RGB}{97,106,123}
\definecolor{accent}{RGB}{199,63,48}
\definecolor{blue}{RGB}{43,96,151}
\definecolor{green}{RGB}{63,125,111}
\definecolor{gold}{RGB}{180,130,35}
\definecolor{line}{RGB}{219,224,232}
\definecolor{softblue}{RGB}{232,239,248}
\definecolor{softred}{RGB}{250,236,232}
\definecolor{softgreen}{RGB}{233,244,240}
\definecolor{softgold}{RGB}{248,240,219}

\setbeamercolor{normal text}{fg=ink,bg=paper}
\setbeamercolor{frametitle}{fg=ink,bg=paper}
\setbeamercolor{structure}{fg=blue}
\setbeamertemplate{navigation symbols}{}
\setbeamertemplate{headline}{}
\setbeamersize{text margin left=0.78cm,text margin right=0.78cm}
\setbeamerfont{frametitle}{series=\bfseries,size=\Large}
\setbeamertemplate{itemize item}{\color{accent}\scriptsize$\blacksquare$}

\setbeamertemplate{footline}{
  \leavevmode
  \hbox{
    \begin{beamercolorbox}[wd=.78\paperwidth,ht=3.0ex,dp=1.2ex,leftskip=.8cm]{author in head/foot}
      \color{muted}\small Pokemon-Datenbankprojekt
    \end{beamercolorbox}
    \begin{beamercolorbox}[wd=.22\paperwidth,ht=3.0ex,dp=1.2ex,rightskip=.8cm]{date in head/foot}
      \color{muted}\small \hfill \insertframenumber/\inserttotalframenumber
    \end{beamercolorbox}
  }
}

\newtcolorbox{panel}[1][]{
  enhanced,
  colback=white,
  colframe=line,
  boxrule=0.7pt,
  arc=3pt,
  left=8pt,
  right=8pt,
  top=7pt,
  bottom=7pt,
  #1
}

\newcommand{\tagbox}[3]{%
  \begin{panel}[colback=#3,colframe=#2!45,height=1.55cm]
    {\Large\bfseries #1}\par
    {\small\color{muted}#2}
  \end{panel}
}

\newcolumntype{Y}{>{\raggedright\arraybackslash}X}
\newcolumntype{L}[1]{>{\raggedright\arraybackslash}p{#1}}

\lstdefinelanguage{SQL}{
  morekeywords={SELECT,FROM,WHERE,JOIN,LEFT,ON,AS,ORDER,BY,GROUP,HAVING,COUNT,AVG,ROUND,IS,NOT,NULL,DESC,ASC,PRAGMA},
  sensitive=false,
  morecomment=[l]{--},
  morestring=[b]'
}

\lstset{
  language=SQL,
  basicstyle=\ttfamily\scriptsize,
  keywordstyle=\color{blue}\bfseries,
  commentstyle=\color{muted},
  stringstyle=\color{accent},
  showstringspaces=false,
  columns=fullflexible,
  keepspaces=true,
  breaklines=true,
  frame=none
}

\title{Pokemon-Datenbankprojekt}
\subtitle{Planung, Modellierung, SQL und Reflexion}
\author{Robert, Joel, Berkay, Felix, David}
\date{}

\begin{document}

\begin{frame}[plain]
  \vspace*{0.9cm}
  {\fontsize{30}{34}\selectfont\bfseries Pokemon-Datenbankprojekt}\par
  \vspace{0.35cm}
  {\Large Planung, Datenmodell, SQL-Abfragen und Reflexion}\par
  \vspace{0.8cm}
  \begin{panel}[colback=softblue,colframe=blue!35]
    \begin{minipage}[t]{0.31\textwidth}
      {\bfseries Fokus}\par
      \small Relationale Struktur statt Einzeltabelle
    \end{minipage}\hfill
    \begin{minipage}[t]{0.31\textwidth}
      {\bfseries Datenbasis}\par
      \small 21 Karten, 29 Attacken-Verknuepfungen
    \end{minipage}\hfill
    \begin{minipage}[t]{0.31\textwidth}
      {\bfseries Ziel}\par
      \small Abfragen erklaeren, Modellentscheidungen begruenden
    \end{minipage}
  \end{panel}
  \vfill
  {\small Team: Robert, Joel, Berkay, Felix, David}
\end{frame}

\begin{frame}[shrink=8]{Agenda}
  \begin{panel}[colback=white]
    \begin{tabularx}{\linewidth}{@{}L{0.08\linewidth}Y Y@{}}
      \toprule
      \textbf{Teil} & \textbf{Kapitel} & \textbf{Leitfrage} \\
      \midrule
      1 & Projektplanung & Was war MUSS/SOLL/KANN und welche Entscheidung hielt das Projekt im Zeitrahmen? \\
      2 & Datenmodell & Welche Entit\"aten, Beziehungen, Schl\"ussel und Normalisierungsideen stecken im Modell? \\
      3 & SQL & Welche JOIN-, GROUP-BY- und HAVING-Abfragen zeigen, dass das Modell funktioniert? \\
      4 & Qualit\"at und Reflexion & Welche Tests, L\"ucken und Verbesserungen ergeben sich f\"ur eine zweite Version? \\
      \bottomrule
    \end{tabularx}
  \end{panel}
\end{frame}

\begin{frame}{Projektplanung: MUSS, SOLL, KANN}
  \begin{minipage}[t]{0.32\textwidth}
    \begin{panel}[colback=softred,colframe=accent!40,height=5.8cm]
      {\Large\bfseries MUSS}\par
      \vspace{0.4em}
      \begin{itemize}
        \item Name
        \item Kartennummer
        \item Typ
        \item HP
        \item Attacken
      \end{itemize}
      \vspace{0.25em}
      \small\color{muted}Ohne diese Felder w\"are die Datenbank fachlich nicht vorzeigbar.
    \end{panel}
  \end{minipage}\hfill
  \begin{minipage}[t]{0.32\textwidth}
    \begin{panel}[colback=softblue,colframe=blue!40,height=5.8cm]
      {\Large\bfseries SOLL}\par
      \vspace{0.4em}
      \begin{itemize}
        \item Illustrator
        \item Rarit\"at
        \item Evolution
      \end{itemize}
      \vspace{0.25em}
      \small\color{muted}Diese Felder machen das Modell interessanter und zeigen echte Beziehungen.
    \end{panel}
  \end{minipage}\hfill
  \begin{minipage}[t]{0.32\textwidth}
    \begin{panel}[colback=softgold,colframe=gold!55,height=5.8cm]
      {\Large\bfseries KANN}\par
      \vspace{0.4em}
      \begin{itemize}
        \item Preis
        \item Schw\"ache
        \item Resistenz
        \item R\"uckzugskosten
      \end{itemize}
      \vspace{0.25em}
      \small\color{muted}Diese Daten sind sinnvoll, aber nicht entscheidend f\"ur den Prototyp.
    \end{panel}
  \end{minipage}
\end{frame}

\begin{frame}{Projektplanung: wichtigste Entscheidung}
  \begin{panel}[colback=softblue,colframe=blue!35]
    {\Large\bfseries Erst Pflichtdaten stabil, dann Erweiterungen.}
  \end{panel}
  \vspace{0.45cm}
  \begin{minipage}[t]{0.48\textwidth}
    \begin{panel}
      {\bfseries Warum das geholfen hat}\par
      \begin{itemize}
        \item Die Datenbank war fr\"uh abfragbar.
        \item Fehler im Modell wurden vor der Abgabe sichtbar.
        \item Zusatzfelder konnten weggelassen werden, ohne das Kernziel zu gef\"ahrden.
      \end{itemize}
    \end{panel}
  \end{minipage}\hfill
  \begin{minipage}[t]{0.48\textwidth}
    \begin{panel}
      {\bfseries Konkrete Folge}\par
      \begin{itemize}
        \item 21 Karten haben Name, Nummer, Typ und HP vollst\"andig.
        \item Attacken sind nicht nur Text, sondern separat modelliert.
        \item Preis und Resistenz wurden als ehrliche L\"ucken dokumentiert.
      \end{itemize}
    \end{panel}
  \end{minipage}
\end{frame}

\begin{frame}[shrink=9]{Datenmodellierung: Entit\"aten}
  \begin{panel}
    \small
    \begin{tabularx}{\linewidth}{@{}L{0.18\linewidth}Y Y@{}}
      \toprule
      \textbf{Entit\"at} & \textbf{Was wird gespeichert?} & \textbf{Warum eigene Tabelle?} \\
      \midrule
      \texttt{card} & Nummer, Name, HP, Verweise & zentrale Entit\"at, jede Auswertung startet hier \\
      \texttt{type} & Feuer, Wasser, Kampf usw. & Typnamen wiederholen sich bei vielen Karten \\
      \texttt{rarity} & Common, Uncommon, Rare usw. & verhindert Schreibvarianten und Redundanz \\
      \texttt{illustrator} & Name des Illustrators & Namen sollen nicht mehrfach als Freitext stehen \\
      \texttt{attack} & Name, Schaden, Beschreibung & Attacken sind eigene Informationen \\
      \texttt{card\_attack} & Karte, Attacke, Reihenfolge, Kosten & technische Umsetzung der mehrwertigen Attackenbeziehung \\
      \bottomrule
    \end{tabularx}
  \end{panel}
  \vspace{0.12cm}
  \scriptsize\color{muted}Warum nicht eine Tabelle? Typen, Rarit\"aten, Illustratoren und Attackentexte w\"urden wiederholt und Korrekturen w\"aren fehleranf\"allig.
\end{frame}

\begin{frame}{Datenmodellierung: Tabellen\"ubersicht statt Diagramm}
  \begin{minipage}[t]{0.49\textwidth}
    \begin{panel}[height=5.9cm]
      {\bfseries Kerntabellen}\par
      \vspace{0.25em}
      \begin{tabularx}{\linewidth}{@{}L{0.32\linewidth}Y@{}}
        \toprule
        \textbf{Tabelle} & \textbf{wichtige Felder} \\
        \midrule
        \texttt{card} & \texttt{id}, \texttt{card\_number}, \texttt{name}, \texttt{hp}, FK-Felder \\
        \texttt{attack} & \texttt{id}, \texttt{name}, \texttt{damage}, \texttt{description} \\
        \texttt{card\_attack} & \texttt{card\_id}, \texttt{attack\_id}, \texttt{attack\_order}, \texttt{energy\_cost} \\
        \bottomrule
      \end{tabularx}
    \end{panel}
  \end{minipage}\hfill
  \begin{minipage}[t]{0.49\textwidth}
    \begin{panel}[height=5.9cm]
      {\bfseries Nachschlagetabellen}\par
      \vspace{0.25em}
      \begin{tabularx}{\linewidth}{@{}L{0.36\linewidth}Y@{}}
        \toprule
        \textbf{Tabelle} & \textbf{Rolle im Modell} \\
        \midrule
        \texttt{type} & wiederverwendbare Typnamen \\
        \texttt{rarity} & wiederverwendbare Seltenheitsstufen \\
        \texttt{illustrator} & wiederverwendbare Namen \\
        \texttt{card} auf \texttt{card} & Evolution als Selbstbezug \\
        \bottomrule
      \end{tabularx}
    \end{panel}
  \end{minipage}
\end{frame}

\begin{frame}[shrink=12]{Kardinalit\"aten: Beziehungen im Modell}
  \begin{panel}
    \small
    \begin{tabularx}{\linewidth}{@{}L{0.20\linewidth}L{0.12\linewidth}Y Y@{}}
      \toprule
      \textbf{Beziehung} & \textbf{Kardinalit\"at} & \textbf{Bedeutung} & \textbf{technische Umsetzung} \\
      \midrule
      Typ zu Karte & 1:n & Ein Typ kann bei vielen Karten vorkommen. & \texttt{card.type\_id} zu \texttt{type.id}. \\
      Rarit\"at zu Karte & 1:n & Eine Seltenheit gilt f\"ur viele Karten. & \texttt{card.rarity\_id} zu \texttt{rarity.id}. \\
      Karte zu Attacke & n:m & Eine Karte kann mehrere Attacken haben. & \texttt{card\_attack} mit \texttt{card\_id} und \texttt{attack\_id}. \\
      Karte zu Vorentwicklung & Selbstbezug & Eine Karte kann sich aus einer anderen Karte entwickeln. & \texttt{evolves\_from\_card\_id} zu \texttt{card.id}. \\
      \bottomrule
    \end{tabularx}
  \end{panel}
\end{frame}

\begin{frame}{Normalisierung: Redundanz vermeiden}
  \begin{minipage}[t]{0.49\textwidth}
    \begin{panel}[colback=softgreen,colframe=green!40,height=5.85cm]
      {\bfseries Beispiel: Rarit\"at und Typ}\par
      \begin{itemize}
        \item \texttt{Common}, \texttt{Uncommon}, \texttt{Rare} stehen in \texttt{rarity}.
        \item \texttt{Feuer}, \texttt{Wasser}, \texttt{Kampf} stehen in \texttt{type}.
        \item In \texttt{card} steht nur die jeweilige ID.
      \end{itemize}
    \end{panel}
  \end{minipage}\hfill
  \begin{minipage}[t]{0.49\textwidth}
    \begin{panel}[colback=softred,colframe=accent!40,height=5.85cm]
      {\bfseries Problem ohne Normalisierung}\par
      \begin{itemize}
        \item Tippfehler wie \texttt{Uncommen} w\"urden eigene Gruppen erzeugen.
        \item \"Anderungen m\"ussten in vielen Kartenzeilen korrigiert werden.
        \item Statistikabfragen w\"urden unzuverl\"assig.
        \item Attacken als Spalten w\"aren unflexibel bei mehr oder weniger Attacken.
      \end{itemize}
    \end{panel}
  \end{minipage}
\end{frame}

\begin{frame}[shrink=8]{Schl\"ussel: Prim\"arschl\"ussel und Fremdschl\"ussel}
  \begin{panel}
    \begin{tabularx}{\linewidth}{@{}L{0.22\linewidth}Y Y@{}}
      \toprule
      \textbf{Kategorie} & \textbf{bei uns} & \textbf{Funktion} \\
      \midrule
      Prim\"arschl\"ussel & \texttt{id} in \texttt{card}, \texttt{type}, \texttt{rarity}, \texttt{illustrator}, \texttt{attack} & identifiziert jeden Datensatz eindeutig \\
      logische Fremdschl\"ussel & z.B. \texttt{card.type\_id}, \texttt{card.rarity\_id}, \texttt{card.illustrator\_id} & verbinden Karte mit Nachschlagetabellen \\
      Zwischenschl\"ussel & \texttt{card\_attack.card\_id} und \texttt{card\_attack.attack\_id} & verbinden Karten und Attacken \\
      Selbstbezug & \texttt{card.evolves\_from\_card\_id} & bildet Evolutionsketten innerhalb derselben Tabelle ab \\
      \bottomrule
    \end{tabularx}
  \end{panel}
  \vspace{0.18cm}
  \small\color{muted}Aktueller Stand: IDs sind vorhanden und inhaltlich konsistent. Echte \texttt{FOREIGN KEY}-Constraints fehlen noch.
\end{frame}

\begin{frame}[shrink=8]{Datenqualit\"at: Constraints und Wirkung}
  \begin{panel}
    \begin{tabularx}{\linewidth}{@{}L{0.18\linewidth}L{0.22\linewidth}Y@{}}
      \toprule
      \textbf{Constraint} & \textbf{aktueller Stand} & \textbf{welchen Fehler er verhindern w\"urde} \\
      \midrule
      \texttt{PRIMARY KEY} & vorhanden f\"ur zentrale Tabellen & doppelte oder nicht eindeutig adressierbare Datens\"atze \\
      \texttt{FOREIGN KEY} & logisch genutzt, technisch nicht angelegt & Verweise auf nicht existierende Typen, Attacken oder Karten \\
      \texttt{NOT NULL} & kaum konsequent gesetzt & Karten ohne Name, Typ, HP oder Nummer \\
      \texttt{UNIQUE} & nicht konsequent gesetzt & doppelte Kartennummern oder doppelte Typnamen \\
      \texttt{CHECK} & nicht genutzt & unplausible Werte wie negative HP oder leere Energiekosten \\
      \bottomrule
    \end{tabularx}
  \end{panel}
\end{frame}

\begin{frame}{SQL-Abfragen: roter Faden}
  \begin{panel}[colback=softblue,colframe=blue!35]
    {\Large\bfseries Jede Abfrage zeigt einen anderen Nutzen des Datenmodells.}
  \end{panel}
  \vspace{0.35cm}
  \begin{minipage}[t]{0.24\textwidth}
    \begin{panel}[height=4.3cm]
      {\bfseries 1 JOIN}\par
      \small Kartenansicht aus mehreren Tabellen zusammensetzen.
    \end{panel}
  \end{minipage}\hfill
  \begin{minipage}[t]{0.24\textwidth}
    \begin{panel}[height=4.3cm]
      {\bfseries 2 GROUP BY}\par
      \small Typen z\"ahlen und Durchschnittswerte bilden.
    \end{panel}
  \end{minipage}\hfill
  \begin{minipage}[t]{0.24\textwidth}
    \begin{panel}[height=4.3cm]
      {\bfseries 3 HAVING}\par
      \small Gruppen nach berechneten Ergebnissen filtern.
    \end{panel}
  \end{minipage}\hfill
  \begin{minipage}[t]{0.24\textwidth}
    \begin{panel}[height=4.3cm]
      {\bfseries 4 SELF JOIN}\par
      \small Evolutionen innerhalb der Kartentabelle auswerten.
    \end{panel}
  \end{minipage}
\end{frame}

\begin{frame}[fragile,shrink=5]{Abfrage 1: Gesamtansicht einer Karte}
  \begin{minipage}[t]{0.55\textwidth}
    \begin{panel}
\begin{lstlisting}
SELECT
  c.card_number,
  c.name,
  t.name AS type,
  c.hp,
  i.name AS illustrator,
  r.name AS rarity
FROM card c
JOIN type t ON t.id = c.type_id
JOIN illustrator i ON i.id = c.illustrator_id
JOIN rarity r ON r.id = c.rarity_id
ORDER BY c.id;
\end{lstlisting}
    \end{panel}
  \end{minipage}\hfill
  \begin{minipage}[t]{0.41\textwidth}
    \begin{panel}[colback=softblue,colframe=blue!35]
      {\bfseries Ergebnis}\par
      \begin{itemize}
        \item 21 lesbare Kartenzeilen
        \item Pflichtdaten plus Typ, Illustrator und Rarit\"at
        \item Beispiel: \textbf{Camerupt}, Feuer, 140 HP, Uncommon
      \end{itemize}
      \vspace{0.2em}
      \small\color{muted}Die ausgelagerten Tabellen lassen sich wieder zu einer kompakten Kartenansicht verbinden.
    \end{panel}
  \end{minipage}
\end{frame}

\begin{frame}{Abfrage 1: Warum genau diese JOINs?}
  \begin{panel}
    \begin{tabularx}{\linewidth}{@{}L{0.27\linewidth}Y Y@{}}
      \toprule
      \textbf{Join-Bedingung} & \textbf{Bedeutung} & \textbf{warum nicht anders?} \\
      \midrule
      \texttt{t.id = c.type\_id} & verbindet Karte mit kontrolliertem Typdatensatz & Textvergleich w\"are anf\"allig f\"ur Schreibfehler \\
      \texttt{i.id = c.illustrator\_id} & holt den Namen aus \texttt{illustrator} & Illustratorname m\"usste sonst in jeder Karte stehen \\
      \texttt{r.id = c.rarity\_id} & holt die Seltenheitsstufe aus \texttt{rarity} & Freitext-Rarit\"aten w\"urden Statistiken verf\"alschen \\
      \bottomrule
    \end{tabularx}
  \end{panel}
  \vspace{0.35cm}
  \begin{panel}[colback=softgold,colframe=gold!45]
    {\bfseries \"Uberleitung:} Wenn eine Einzelkarte sauber rekonstruierbar ist, kann man im n\"achsten Schritt alle Karten zusammen statistisch auswerten.
  \end{panel}
\end{frame}

\begin{frame}[fragile]{Abfrage 2: Typverteilung mit GROUP BY}
  \begin{minipage}[t]{0.50\textwidth}
    \begin{panel}
\begin{lstlisting}
SELECT
  t.name AS type,
  COUNT(*) AS cards,
  ROUND(AVG(c.hp), 0) AS avg_hp
FROM card c
JOIN type t ON t.id = c.type_id
GROUP BY t.name
ORDER BY cards DESC, t.name;
\end{lstlisting}
    \end{panel}
  \end{minipage}\hfill
  \begin{minipage}[t]{0.46\textwidth}
    \begin{panel}[colback=softgold,colframe=gold!45]
      {\bfseries Ergebnis}\par
      \vspace{0.2em}
      \begin{tabularx}{\linewidth}{@{}Yrr@{}}
        \toprule
        Typ & Karten & Ø HP \\
        \midrule
        Feuer & 4 & 165 \\
        Wasser & 4 & 95 \\
        Kampf & 3 & 77 \\
        Pflanze & 3 & 80 \\
        Psycho & 3 & 97 \\
        Elektro & 2 & 140 \\
        Finsternis & 2 & 70 \\
        \bottomrule
      \end{tabularx}
    \end{panel}
  \end{minipage}
\end{frame}

\begin{frame}{Abfrage 2: Aussage der Statistik}
  \begin{minipage}[t]{0.48\textwidth}
    \begin{panel}[colback=softgreen,colframe=green!40,height=5.4cm]
      {\bfseries Was gewinnt man daraus?}\par
      \begin{itemize}
        \item Feuer und Wasser sind am h\"aufigsten vertreten.
        \item Feuer hat im Datensatz den h\"ochsten durchschnittlichen HP-Wert.
        \item Die Typ-Tabelle macht Gruppierung sauber, weil Typen nicht als Freitext verstreut sind.
      \end{itemize}
    \end{panel}
  \end{minipage}\hfill
  \begin{minipage}[t]{0.48\textwidth}
    \begin{panel}[colback=softblue,colframe=blue!35,height=5.4cm]
      {\bfseries \"Uberleitung}\par
      \begin{itemize}
        \item \texttt{GROUP BY} fasst Datens\"atze zusammen.
        \item Danach kann man mit \texttt{HAVING} genau diese Gruppen filtern.
        \item Das ist passend f\"ur Fragen wie: Welche Karten haben genau zwei Attacken?
      \end{itemize}
    \end{panel}
  \end{minipage}
\end{frame}

\begin{frame}[fragile]{Abfrage 3: Karten mit genau zwei Attacken}
  \begin{minipage}[t]{0.50\textwidth}
    \begin{panel}
\begin{lstlisting}
SELECT
  c.name,
  COUNT(*) AS attacks,
  c.hp
FROM card c
JOIN card_attack ca ON ca.card_id = c.id
GROUP BY c.id
HAVING COUNT(*) = 2
ORDER BY c.hp DESC, c.name;
\end{lstlisting}
    \end{panel}
  \end{minipage}\hfill
  \begin{minipage}[t]{0.46\textwidth}
    \begin{panel}[colback=softgreen,colframe=green!40]
      {\bfseries Ergebnis: 8 Treffer}\par
      \vspace{0.2em}
      \begin{tabularx}{\linewidth}{@{}YY@{}}
        Ns Gelatwino & Ns Gelatroppo \\
        Camerupt & Erikas Ultrigaria \\
        Team Rockets Kokowei & Banette \\
        Team Rockets Digdri & Ns Flampion \\
      \end{tabularx}
      \vspace{0.35em}
      \small\color{muted}\texttt{HAVING} filtert erst nach dem Z\"ahlen. Genau daf\"ur braucht man die Gruppierung pro Karte.
    \end{panel}
  \end{minipage}
\end{frame}

\begin{frame}{Abfrage 3: Warum das Modell hier hilft}
  \begin{panel}[colback=softgold,colframe=gold!45]
    {\Large\bfseries Die Zwischentabelle macht Attacken flexibel.}
  \end{panel}
  \vspace{0.35cm}
  \begin{minipage}[t]{0.48\textwidth}
    \begin{panel}
      {\bfseries Mit \texttt{card\_attack}}\par
      \begin{itemize}
        \item beliebig viele Attacken pro Karte m\"oglich
        \item Z\"ahlen mit \texttt{COUNT(*)}
        \item Filtern mit \texttt{HAVING}
      \end{itemize}
    \end{panel}
  \end{minipage}\hfill
  \begin{minipage}[t]{0.48\textwidth}
    \begin{panel}
      {\bfseries Ohne Zwischentabelle}\par
      \begin{itemize}
        \item feste Spalten wie \texttt{attack1}, \texttt{attack2}
        \item schwieriger bei drei oder mehr Attacken
        \item schlechtere Statistikabfragen
      \end{itemize}
    \end{panel}
  \end{minipage}
\end{frame}

\begin{frame}[fragile]{Abfrage 4: Evolutionen per Self Join}
  \begin{minipage}[t]{0.50\textwidth}
    \begin{panel}
\begin{lstlisting}
SELECT
  c.name AS card,
  e.name AS evolves_from
FROM card c
LEFT JOIN card e
  ON e.id = c.evolves_from_card_id
WHERE c.evolves_from_card_id IS NOT NULL
ORDER BY c.id;
\end{lstlisting}
    \end{panel}
  \end{minipage}\hfill
  \begin{minipage}[t]{0.46\textwidth}
    \begin{panel}[colback=softblue,colframe=blue!35]
      {\bfseries Ergebnis}\par
      \vspace{0.2em}
      \begin{tabularx}{\linewidth}{@{}YY@{}}
        \toprule
        Karte & entwickelt aus \\
        \midrule
        Digdri & Digda \\
        Camerupt & Camaub \\
        Gelatroppo & Gelatini \\
        Gelatwino & Gelatroppo \\
        Kokowei & Owei \\
        \bottomrule
      \end{tabularx}
    \end{panel}
  \end{minipage}
\end{frame}

\begin{frame}{Abfrage 4: Aussage und \"Uberleitung}
  \begin{minipage}[t]{0.48\textwidth}
    \begin{panel}[colback=softgreen,colframe=green!40,height=5.2cm]
      {\bfseries Aussage}\par
      \begin{itemize}
        \item 5 echte Evolutions-Links sind verkn\"upft.
        \item Die Tabelle \texttt{card} wird mit sich selbst verbunden.
        \item Dadurch braucht man keine eigene Evolutionstabelle f\"ur diesen einfachen Fall.
      \end{itemize}
    \end{panel}
  \end{minipage}\hfill
  \begin{minipage}[t]{0.48\textwidth}
    \begin{panel}[colback=softred,colframe=accent!40,height=5.2cm]
      {\bfseries Grenze der Daten}\par
      \begin{itemize}
        \item Im Rohdatensatz gibt es 9 Entwicklungskarten.
        \item Nur 5 konnten sauber auf vorhandene Karten zeigen.
        \item Das leitet direkt zu Datenqualit\"at und Tests \"uber.
      \end{itemize}
    \end{panel}
  \end{minipage}
\end{frame}

\begin{frame}[fragile]{Testen und Debugging}
  \begin{minipage}[t]{0.49\textwidth}
    \begin{panel}
      {\bfseries Smoke Tests}\par
      \begin{itemize}
        \item Anzahl der Karten: 21
        \item Anzahl der Attacken-Verkn\"upfungen: 29
        \item Pflichtfelder leer? Nein
        \item doppelte Kartennummern? Nein
      \end{itemize}
    \end{panel}
  \end{minipage}\hfill
  \begin{minipage}[t]{0.49\textwidth}
    \begin{panel}
      {\bfseries Beispiel-Pr\"ufabfragen}\par
\begin{lstlisting}
SELECT COUNT(*) FROM card;
SELECT COUNT(*) FROM card_attack;
SELECT card_number
FROM card
GROUP BY card_number
HAVING COUNT(*) > 1;
PRAGMA foreign_key_list(card);
\end{lstlisting}
    \end{panel}
  \end{minipage}
  \vspace{0.25cm}
  \small\color{muted}Ergebnis: Die Inhalte passen logisch zusammen, aber \texttt{PRAGMA foreign\_key\_list} zeigt keine technisch angelegten Fremdschl\"ussel.
\end{frame}

\begin{frame}{Datenl\"ucken: Was fehlt f\"ur eine bessere DB?}
  \begin{panel}
    \begin{tabularx}{\linewidth}{@{}L{0.24\linewidth}L{0.18\linewidth}Y@{}}
      \toprule
      \textbf{Datenbereich} & \textbf{Stand} & \textbf{Bedeutung} \\
      \midrule
      Pflichtfelder & vollst\"andig & Name, Nummer, Typ und HP sind bei allen 21 Karten vorhanden. \\
      Resistenz & 1 von 21 & fachlich fast leer, deshalb nur vorsichtig auswertbar \\
      Preis & 0 von 21 & als KANN-Anforderung nicht umgesetzt \\
      Evolutionen & 5 echte Links bei 9 Entwicklungskarten & mehrere Vorstufen fehlen in der Auswahl \\
      Constraints & teilweise nur logisch & Datenbank verhindert falsche Eingaben noch nicht selbst \\
      \bottomrule
    \end{tabularx}
  \end{panel}
\end{frame}

\begin{frame}{Softwareentwicklung und Reflexion}
  \begin{minipage}[t]{0.49\textwidth}
    \begin{panel}[colback=softred,colframe=accent!40,height=5.8cm]
      {\bfseries Was w\"urden wir \"andern?}\par
      \begin{itemize}
        \item echte \texttt{FOREIGN KEY}-Constraints im Schema
        \item \texttt{NOT NULL}, \texttt{UNIQUE}, \texttt{CHECK} konsequent setzen
        \item \texttt{card\_new} entfernen oder sinnvoll integrieren
        \item Importskript statt manueller Einpflege
      \end{itemize}
    \end{panel}
  \end{minipage}\hfill
  \begin{minipage}[t]{0.49\textwidth}
    \begin{panel}[colback=softblue,colframe=blue!40,height=5.8cm]
      {\bfseries N\"achster sinnvoller Schritt}\par
      \begin{itemize}
        \item Set-Tabelle mit Setname und Setnummer
        \item Preis als eigenes Feld oder Verlaufstabelle
        \item vollst\"andige Evolutionsketten importieren
        \item kleine Such- oder Filteroberfl\"ache f\"ur Karten
      \end{itemize}
    \end{panel}
  \end{minipage}
\end{frame}

\begin{frame}{Fazit}
  \begin{panel}[colback=softgreen,colframe=green!40]
    {\Large\bfseries Die Grundstruktur funktioniert.}\par
    \vspace{0.3em}
    \large Pflichtdaten sind vollst\"andig, zentrale Beziehungen sind nachvollziehbar modelliert und SQL-Abfragen zeigen die Vorteile der Aufteilung.
  \end{panel}
  \vspace{0.4cm}
  \begin{panel}[colback=softgold,colframe=gold!45]
    {\Large\bfseries Die zweite Version w\"are technischer.}\par
    \vspace{0.3em}
    \large Mehr Constraints, automatischer Import, bessere Zusatzdaten und vollst\"andige Evolutionsketten.
  \end{panel}
  \vfill
  \centering{\Large\bfseries Fragen?}
\end{frame}

\end{document}
